% Clase 23 --- Diagrama fasorial de impedancias en el plano complejo (§1.8).
% El docente lo dibuja para leer |Z| y cos(phi) de la geometria, sin cuentas:
% R sobre el eje real, la bobina hacia arriba, el condensador hacia abajo.
% Dos paneles con la MISMA geometria; lo unico que cambia es cual de las dos
% reactancias gana.
\begin{center}
\begin{tikzpicture}[scale=1]

% ---------------- (a) circuito inductivo: omega L > 1/(omega C) ------------
\begin{scope}
  \draw[figaxis] (-0.4,0) -- (4.7,0) node[figlblaux, anchor=south east] {Re};
  \draw[figaxis] (0,-2.0) -- (0,2.7) node[figlblaux, anchor=south east] {Im};

  % R sobre el eje real
  \draw[figvec] (0,0) -- (3,0);
  \node[figlbl, anchor=north] at (1.5,-0.12) {$R$};

  % bobina hacia arriba y condensador hacia abajo, encadenados en la punta de R
  \draw[figvec] (3,0) -- (3,2.2);
  \node[figlbl, anchor=west] at (3.08,0.6) {$\omega L$};
  \draw[figvec] (3.6,2.2) -- (3.6,1.3);
  \node[figlbl, anchor=west] at (3.7,1.78) {$-\dfrac{1}{\omega C}$};
  \draw[figaux] (3,2.2) -- (3.6,2.2);
  \draw[figaux] (3,1.3) -- (3.6,1.3);

  % impedancia total
  \draw[figvecres] (0,0) -- (3,1.3);
  \node[figlbl, figred, anchor=south east] at (2.5,1.09) {$\hat Z$};

  \draw[figgray, line width=0.6pt] (0.95,0) arc (0:23.4:0.95);
  \node[figlbl, anchor=west] at (1.0,0.26) {$\phi > 0$};

  \node[figlbl, anchor=north, align=center] at (2.0,-2.35)
    {\textbf{(a)} $\omega L > \dfrac{1}{\omega C}$: inductivo,\\
     la corriente atrasa};
\end{scope}

% ---------------- (b) circuito capacitivo: omega L < 1/(omega C) -----------
\begin{scope}[xshift=6.4cm]
  \draw[figaxis] (-0.4,0) -- (4.7,0) node[figlblaux, anchor=south east] {Re};
  \draw[figaxis] (0,-2.0) -- (0,2.7) node[figlblaux, anchor=south east] {Im};

  \draw[figvec] (0,0) -- (3,0);
  \node[figlbl, anchor=south] at (1.5,0.12) {$R$};

  \draw[figvec] (3,0) -- (3,0.9);
  \node[figlbl, anchor=east] at (2.9,0.62) {$\omega L$};
  \draw[figvec] (3.6,0.9) -- (3.6,-1.3);
  \node[figlbl, anchor=east] at (3.5,-0.5) {$-\dfrac{1}{\omega C}$};
  \draw[figaux] (3,0.9) -- (3.6,0.9);
  \draw[figaux] (3,-1.3) -- (3.6,-1.3);

  \draw[figvecres] (0,0) -- (3,-1.3);
  \node[figlbl, figred, anchor=north east] at (2.4,-1.0) {$\hat Z$};

  \draw[figgray, line width=0.6pt] (0.95,0) arc (0:-23.4:0.95);
  \node[figlbl, anchor=west] at (1.0,-0.26) {$\phi < 0$};

  \node[figlbl, anchor=north, align=center] at (2.0,-2.35)
    {\textbf{(b)} $\omega L < \dfrac{1}{\omega C}$: capacitivo,\\
     la corriente adelanta};
\end{scope}

\end{tikzpicture}\\[0.5em]
{\small\itshape Las tres impedancias no se suman como números sino como
vectores del plano complejo, y ahí está toda la gracia: la resistencia va sobre
el eje real, la bobina hacia arriba y el condensador hacia abajo, de modo que
las dos reactancias \textbf{se restan entre sí} antes de componerse con $R$. El
triángulo rectángulo que queda dice de un vistazo lo que costaría media carilla
de trigonometría: la hipotenusa es $|Z|$, el cateto adyacente es $R$, y por lo
tanto el factor de potencia es $\cos\phi = R/|Z|$. Como sólo la componente
\emph{real} sobrevive en $\langle P\rangle = \tfrac12\mathcal{E}_0 I_m\cos\phi$,
el dibujo también explica por qué únicamente la resistencia disipa: si el
triángulo es todo cateto vertical ($R = 0$), el vector $\hat Z$ es puramente
imaginario, $\phi = \pm\pi/2$ y la potencia media es cero. En resonancia
($\omega L = 1/\omega C$) los dos aportes verticales se cancelan, $\hat Z$
queda acostado sobre el eje real y el factor de potencia vale $1$.}
\end{center}
